% Cong thuc xep hang rui ro lo hong (Risk Score)
% Bien dich bang XeLaTeX:  xelatex cong_thuc_risk_score.tex
\documentclass[12pt,a4paper]{article}

\usepackage{fontspec}
\setmainfont{Times New Roman}
\usepackage{amsmath,amssymb}
\usepackage[margin=2.2cm]{geometry}
\usepackage{booktabs}
\usepackage{array}
\usepackage{enumitem}
\usepackage[hidelinks]{hyperref}

\renewcommand{\refname}{Tài liệu tham khảo}

\setlength{\parindent}{0pt}
\setlength{\parskip}{6pt}
\renewcommand{\arraystretch}{1.3}

\begin{document}

\begin{center}
  {\Large\bfseries Công thức xếp hạng rủi ro lỗ hổng (Risk Score)}\\[4pt]
  Học phần Chính sách và Quản trị an ninh mạng, Nhóm 14
\end{center}

\section{Công thức}

Điểm rủi ro nhận giá trị từ 0 đến 100, kết hợp bốn yếu tố: mức quan trọng của tài sản,
danh mục CISA KEV, điểm CVSS và điểm EPSS.

\begin{equation}
\label{eq:risk}
\begin{split}
\text{Risk Score} = {}& \text{Asset Criticality} \times \Bigl[ \bigl( K_{\text{KEV}} \times 10 \bigr) \\
  &+ \bigl( 1 - K_{\text{KEV}} \bigr) \times \bigl( \text{CVSS} \times \text{EPSS}^{0{,}3} \bigr) \Bigr] \times 10
\end{split}
\end{equation}

Vì $K_{\text{KEV}}$ chỉ nhận giá trị 0 hoặc 1, công thức tương đương với dạng hai nhánh:

\begin{equation}
\text{Risk Score} =
\begin{cases}
  100 \times \text{Asset Criticality}, & \text{nếu } K_{\text{KEV}} = 1 \\[6pt]
  10 \times \text{Asset Criticality} \times \text{CVSS} \times \text{EPSS}^{0{,}3}, & \text{nếu } K_{\text{KEV}} = 0
\end{cases}
\end{equation}

\section{Cơ sở của công thức}

Shimizu và Hashimoto~\cite{shimizu2026vmc} đề xuất phương pháp Vulnerability Management Chaining,
xếp ưu tiên lỗ hổng bằng một cây quyết định hai tầng kết hợp KEV, EPSS và CVSS:

\begin{center}
\begin{tabular}{>{\raggedright}p{4.2cm} >{\raggedright}p{6.2cm} >{\raggedright\arraybackslash}p{4.6cm}}
\toprule
\textbf{Tầng} & \textbf{Điều kiện} & \textbf{Nếu không đạt} \\
\midrule
1. Lọc theo mối đe dọa & Có trong KEV, hoặc EPSS từ 0,088 trở lên & Dời sang chu kỳ vá thông thường \\
2. Đánh giá mức nghiêm trọng & CVSS từ 7,0 trở lên & Ưu tiên thấp hơn \\
\bottomrule
\end{tabular}
\end{center}

Trên 28.377 lỗ hổng công bố từ tháng 4 năm 2022 đến tháng 4 năm 2023, cách tiếp cận này cho hiệu quả
gấp 14 đến 18 lần so với cách chỉ dùng CVSS, với độ bao phủ 85,6\% số lỗ hổng bị khai thác, và giảm
khoảng 95\% khối lượng vá khẩn cấp~\cite{shimizu2026vmc}.

Nhóm giữ nguyên nguyên tắc của bài báo là xét KEV trước, rồi mới tới EPSS và CVSS. Nhóm bổ sung yếu tố
mức quan trọng của tài sản, vốn được chính bài báo nêu là chưa xét tới, và biểu diễn lại toàn bộ thành
công thức~\eqref{eq:risk}. Phần bổ sung này, cùng số mũ 0,3 và hệ số nhân 10, là đề xuất của nhóm,
không có trong bài báo gốc.

\section{Giải thích các biến}

\begin{center}
\begin{tabular}{>{\raggedright}p{3.4cm} >{\raggedright}p{3.8cm} >{\raggedright\arraybackslash}p{7.4cm}}
\toprule
\textbf{Biến} & \textbf{Miền giá trị} & \textbf{Ý nghĩa} \\
\midrule
Asset Criticality & $\{0{,}2;\ 0{,}5;\ 0{,}8;\ 1{,}0\}$ & Mức quan trọng của tài sản chứa lỗ hổng \\
$K_{\text{KEV}}$ & $\{0,\ 1\}$ & Bằng 1 nếu lỗ hổng có trong danh mục CISA KEV, tức đã được xác nhận bị khai thác thực tế; bằng 0 nếu không \\
CVSS & $[0,\ 10]$ & Điểm CVSS Base, đo mức nghiêm trọng kỹ thuật \\
EPSS & $[0,\ 1]$ & Xác suất lỗ hổng bị khai thác trong 30 ngày tới, do FIRST công bố hằng ngày \\
$\text{EPSS}^{0{,}3}$ & $[0,\ 1]$ & Nâng các giá trị EPSS nhỏ lên để không bị điểm CVSS lấn át \\
\bottomrule
\end{tabular}
\end{center}

\section{Quy đổi mức quan trọng của tài sản}

\begin{center}
\begin{tabular}{>{\raggedright}p{2.6cm} >{\centering}p{2.4cm} >{\raggedright\arraybackslash}p{10cm}}
\toprule
\textbf{Mức} & \textbf{Hệ số} & \textbf{Ví dụ} \\
\midrule
Critical & 1,0 & Hệ thống cốt lõi, cơ sở dữ liệu khách hàng, hệ thống mở ra Internet \\
High & 0,8 & Máy chủ nội bộ, Active Directory \\
Medium & 0,5 & Máy trạm nhân viên, máy chủ staging \\
Low & 0,2 & Thiết bị thử nghiệm, máy in, phòng lab \\
\bottomrule
\end{tabular}
\end{center}

\section{Ví dụ tính toán}

\begin{center}
\begin{tabular}{>{\raggedright}p{4.6cm} c c c c >{\raggedright\arraybackslash}p{2.2cm}}
\toprule
\textbf{Kịch bản} & \textbf{AC} & \textbf{KEV} & \textbf{CVSS} & \textbf{EPSS} & \textbf{Risk Score} \\
\midrule
Lỗ hổng đang bị khai thác trên máy chủ cốt lõi & 1,0 & 1 & 10,0 & & 100 \\
Lỗ hổng mới trên máy chủ cốt lõi, dự báo khả năng bị khai thác trung bình & 1,0 & 0 & 7,5 & 0,45 & 59,0 \\
Lỗ hổng đang bị khai thác trên máy tính lễ tân & 0,2 & 1 & 8,0 & & 20,0 \\
Lỗ hổng nặng trên máy thử nghiệm, chưa có dấu hiệu khai thác & 0,2 & 0 & 9,8 & 0,001 & 2,47 \\
\bottomrule
\end{tabular}
\end{center}

Trong bảng, AC là viết tắt của Asset Criticality. Chi tiết phép tính của kịch bản 2 và kịch bản 4:
\begin{align*}
\text{Kịch bản 2:}\quad & 1{,}0 \times 7{,}5 \times 0{,}45^{0{,}3} \times 10 = 1{,}0 \times 7{,}5 \times 0{,}787 \times 10 \approx 59{,}0 \\
\text{Kịch bản 4:}\quad & 0{,}2 \times 9{,}8 \times 0{,}001^{0{,}3} \times 10 = 0{,}2 \times 9{,}8 \times 0{,}126 \times 10 \approx 2{,}47
\end{align*}

\section{Lưu ý khi áp dụng}

\begin{itemize}[leftmargin=1.4em]
  \item Số mũ 0,3 được chọn theo kinh nghiệm, chưa được hiệu chỉnh bằng dữ liệu.
  \item Ở nhánh $K_{\text{KEV}} = 1$, điểm chỉ phụ thuộc vào mức quan trọng của tài sản. Một lỗ hổng
        đang bị khai thác trên tài sản mức Low chỉ được 20 điểm, dù máy bị chiếm có thể trở thành
        bàn đạp để tấn công sang các máy khác. Vì vậy mọi lỗ hổng trong KEV vẫn nên có hạn vá bắt
        buộc, mức quan trọng của tài sản chỉ quyết định hạn vá dài hay ngắn.
  \item Điểm EPSS thay đổi hằng ngày, nên điểm rủi ro cần được tính lại hằng ngày thay vì chỉ tính một lần.
  \item Bài báo của Shimizu và Hashimoto lấy điểm EPSS tại một thời điểm duy nhất vào cuối kỳ nghiên cứu
        cho mọi lỗ hổng. Điểm EPSS lấy muộn có thể đã phản ánh việc lỗ hổng bị khai thác, nên mức hiệu quả
        báo cáo có thể cao hơn khi áp dụng thực tế, lúc doanh nghiệp phải quyết định ngay sau khi lỗ hổng
        được công bố.
\end{itemize}

\begin{thebibliography}{9}
\bibitem{shimizu2026vmc}
N.~Shimizu and M.~Hashimoto, ``Vulnerability Management Chaining: An Integrated Framework for Efficient
Cybersecurity Risk Prioritization,'' \emph{IEEE Access}, vol.~14, pp.~31407--31424, 2026,
doi: \href{https://doi.org/10.1109/ACCESS.2026.3665768}{10.1109/ACCESS.2026.3665768}.
Bản tiền ấn phẩm: \href{https://arxiv.org/abs/2506.01220}{arXiv:2506.01220}.
\end{thebibliography}

\end{document}
